% =====================================================================
% background.tex --- IWAI 2026, Hallgren & Hyland
% Model-walkthrough brief (2026-05-30 reshape).
%
% Structure: a step-by-step tour of the underlying mathematical model.
% Each step opens with the question the previous step leaves
% unanswered, and shows the next formal piece as the natural answer.
%
% Reuse: equations and figures lifted verbatim from
% paper/pomdp_paradigm_model.tex (eq:gen, eq:discrim, eq:vfe,
% eq:motiv, eq:fast, eq:marglik, eq:slow, fig:factor, fig:trust,
% fig:anticipated). Captions adapted to the walkthrough narrative.
%
% Predecessor (lit-review version): background_litreview.tex in the
% same directory, archived for reference. Source-section drafts S{N}_*.tex
% remain in place.
%
% Citation key handling: keys with the hand-curated style
% (e.g. kuhn1962Structure) resolve against paper/refs.bib. Keys of
% the form Author2024_3 are Semantic Scholar auto-keys from the
% cluster bibs at paper/bib_additions/cluster_C*.bib. Foundational
% anchors flagged % TODO ADD need to be hand-added to refs.bib:
%   kunda1990MotivatedReasoning, zollman2007NetworkEpistemology,
%   lakatos1970FalsificationMethodology.
% =====================================================================

% =====================================================================
\section{The Phenomenon}
\label{sec:phenomenon}
% =====================================================================

Scientific communities revise their theoretical commitments more
slowly than the evidence alone would dictate. Anomalies accumulate;
the prevailing framework persists; the shift, when it comes, is
abrupt and path-dependent \citep{kuhn1962Structure}. The interesting
question is not why belief lags --- a Bayesian agent with a sharp
prior is the obvious proximate cause --- but \emph{what generates
the gradient against which belief is moving}, and why that gradient
extends across the community rather than being a quirk of any single
mind \citep{planck1949Autobiography,bourdieu1975ScientificField}.
Kuhn's deepest observation is that the paradigm shapes what the
scientist is even capable of perceiving, of choosing to investigate,
and of entertaining as a live alternative: three interconnected
\emph{faces} of one structural phenomenon. The model we build is
committed to recovering all three from one mechanism, not from three
separately-added biases.

% =====================================================================
\section*{Desiderata}
\addcontentsline{toc}{section}{Desiderata}
\label{sec:desiderata}
% =====================================================================

\noindent\fbox{\parbox{0.97\textwidth}{%
The construction must satisfy the following, in priority order:
\begin{enumerate}\setlength{\itemsep}{2pt}\setlength{\topsep}{2pt}
\item \textbf{One mechanism for three faces.} Theory-laden perception
  (what is seen), theory-laden experimentation (what is investigated),
  and theory-laden hypothesis space (what alternatives are entertained)
  emerge from a single objective.
\item \textbf{Asymmetric, path-dependent revision.} The cost of
  changing one's mind has a principled formal counterpart
  \citep{hylandAlbarracin2025MindChange}, and belief trajectories are
  history-dependent.
\item \textbf{Paradigm-dependent attention as derived behaviour.}
  Which observations the agent samples is conditioned by the current
  paradigm through the policy itself, not a bespoke bias term.
\item \textbf{Network-epistemology phenomena as derived consequences.}
  The speed--accuracy tradeoff of communication graphs,
  paradoxical-stubbornness regularisation, and the
  bubble-versus-chamber distinction \citep{Nguyen2018_1} fall out
  of one substrate, not separate appeals.
\item \textbf{Active-inference machinery, unmodified.} Variational
  and expected free energy \citep{friston2015ActiveInferenceEpistemic,
  millidge2021EFE} enter in standard form; no bespoke extensions.
\item \textbf{Minimality.} The simplest construction satisfying
  (1)--(5).
\end{enumerate}}}

\bigskip\noindent The remainder of this section builds the model
piece by piece. Each step opens with the question its predecessor
leaves open and shows the next formal object as the natural answer.

% =====================================================================
\section{Building the Model}
\label{sec:walkthrough}
% =====================================================================

\subsection{Step 1: Beliefs on a Bayes net}
\label{ssec:bayesnet}

\noindent\textit{What does the agent need to represent?} A paradigm
is not a property of the world but of the agent's reading of the
world: the same physical event yields different evidential
implications under different frameworks \citep{kuhn1962Structure}.
Two latent variables therefore need to be carried, not one. The fast
hidden state $s_t \in \mathcal S$ is the regime realised in an
experiment and carries nontrivial dynamics. The slow latent
$\theory \in \{0, \dots, K{-}1\}$ is the \emph{paradigm} under which
the agent reads $s_t$. The paradigm is not resampled; it is inferred.

A Bayes net falls out of the conditional-independence structure: the
joint factorises as
\begin{equation}
p(o_{1:T}, s_{1:T}, \theory)
\;=\;
p(\theory)\prod_{t} p(o_t \mid s_t, \theory)\,
   p(s_t \mid s_{t-1}, a_{t-1}),
\label{eq:gen}
\end{equation}
with $a_{t-1}$ the agent's action at the previous step. The slow
factor has a near-identity transition $B_\theory \approx I$: the
paradigm does not move on its own. Figure~\ref{fig:factor} shows the
graph; the shape is not a modelling choice but what the independence
structure \emph{requires} once paradigm and world state are accepted
as distinct hidden quantities.

\begin{figure}[t]
\centering
\begin{tikzpicture}[node distance=12mm]
\node[slow] (th) {$\theory$};
\node[lat] (s2) [below=16mm of th] {$s_t$};
\node[lat] (s1) [left=20mm of s2] {$s_{t-1}$};
\node[lat] (s3) [right=20mm of s2] {$s_{t+1}$};
\node[obs] (o1) [below=10mm of s1] {$o_{t-1}$};
\node[obs] (o2) [below=10mm of s2] {$o_{t}$};
\node[obs] (o3) [below=10mm of s3] {$o_{t+1}$};
\draw[->,world] (s1) -- node[above,font=\scriptsize]{$B_s$} (s2);
\draw[->,world] (s2) -- node[above,font=\scriptsize]{$B_s$} (s3);
\draw[->,paradA] (s1) -- (o1);
\draw[->,paradA] (s2) -- node[right,font=\scriptsize]{$p(o\mid s,\theory)$} (o2);
\draw[->,paradA] (s3) -- (o3);
\draw[->,util,dashed] (th) to[out=200,in=70] (o1.north);
\draw[->,util,dashed] (th) -- (o2.north);
\draw[->,util,dashed] (th) to[out=-20,in=110] (o3.north);
\draw[->,util] (th) to[out=150,in=30,looseness=6] node[above,font=\scriptsize]{$B_\theory\!\approx\! I$} (th);
\node[font=\scriptsize, world, align=center] at (s3.east) [right=2mm] {fast\\(world)};
\node[font=\scriptsize, util, align=center] at (th.east) [right=10mm] {slow\\(paradigm)};
\end{tikzpicture}
\caption{The Bayes net for paradigm-dependent belief. The world
state $s_t$ (grey) evolves under nontrivial dynamics $B_s$ and emits
observations through the likelihood $p(o_t\mid s_t,\theory)$ (blue).
The paradigm $\theory$ (orange) is a slow latent
($B_\theory\!\approx\!I$) that conditions every likelihood (dashed):
it does not move the world; it sets how the world is read. The agent
infers both factors.}
\label{fig:factor}
\end{figure}

\subsection{Step 2: The theory-laden likelihood}
\label{ssec:likelihood}

\noindent\textit{What makes $\theory$ a \emph{paradigm} rather than
just another latent?} What separates $\theory$ from $s_t$ in
Eq.~\eqref{eq:gen} is that $\theory$ conditions the
\emph{likelihood} $p(o_t \mid s_t, \theory)$: the same world state
yields different outcome distributions when read through different
paradigms. This is the formal counterpart of theory-laden
perception. The Lavoisian and the phlogistonian were not assigning
different probabilities to a shared proposition; they were operating
with different generative likelihoods \citep{Chang2012_9,
Carrier2002_4}. The residue resisting deflation of incommensurability
is exactly this likelihood-level structural difference, not a
difference in parameter values within one model \citep{Malone1993_5}.

The discriminating capacity of an experiment $a$ is then
\begin{equation}
d(a) \;=\; \tfrac{1}{2}\sum_{\theory \neq \theory'}\!
   \KL{p(o \mid s, \theory)}{p(o \mid s, \theory')},
\label{eq:discrim}
\end{equation}
which is small when paradigms predict similar outcomes (the
experiment is underdetermined) and large when they predict different
ones. The formula records why some experiments do epistemic work in
a paradigm dispute and others do not: only when $d(a)$ is large can
$\theory$ be moved by the result. We make no metaphysical claim that
the world possesses a true paradigm. The world generates observations
through some $\tstar$ which need not lie in the agent's hypothesis
set; the agent maintains a posterior over which $\theory$ best
predicts its experience, and ``which paradigm is true'' reduces to
predictive goodness-of-fit.

\subsection{Step 3: The motivated belief-utility objective}
\label{ssec:motivated}

\noindent\textit{Why doesn't cost-free Bayes already produce paradigm
persistence?} A purely Bayesian agent revises beliefs in proportion
to the strength of evidence; the historical lag of paradigm change
therefore indicates that belief revision carries costs the cost-free
Bayes calculus omits \citep{kunda1990MotivatedReasoning,
hylandAlbarracin2025MindChange}.% TODO ADD kunda1990MotivatedReasoning

The right entry point is the variational free energy of an
observation $o$,
\begin{equation}
F[\qb, o] \;=\;
\underbrace{-\,\E{\log p(o \mid s)}_{\qb}}_{\text{accuracy}}
+\underbrace{\KL{\qb(s)}{p(s)}}_{\text{complexity}},
\label{eq:vfe}
\end{equation}
the sum of an accuracy term and a complexity term measuring
divergence of the posterior from the prior
\citep{friston2015ActiveInferenceEpistemic,mathys2014HGF}.
\citet{hylandAlbarracin2025MindChange} reinterpret the complexity
term of Eq.~\eqref{eq:vfe} as a substantive \emph{cost of changing
one's mind} and recast belief updating as the maximisation of a
belief-utility net of that cost. For a belief $\qb$ over a latent
with carried-over prior $q_0$ and evidence summarised by a marginal
likelihood $L$,
\begin{equation}
q' \;=\;
\arg\max_{\qb}\Bigl\{\,
   U[\qb] + \alpha\,\E{\log L}_{\qb} - \lambda\,\KL{\qb}{q_0}\,
\Bigr\}
\;\;\Longrightarrow\;\;
q' \;\propto\; q_0\, L^{\alpha/\lambda}\, e^{\,u/\lambda},
\label{eq:motiv}
\end{equation}
a tempered Bayesian rule with effective learning rate $\alpha/\lambda$:
$\alpha$ scales the pull of evidence, $\lambda$ the cost of moving,
and the affective utility $u$ (with $U[\qb] = \E{u}_{\qb}$) a
preference over which value is believed, independent of fit. Ordinary
Bayes is the corner $\alpha = \lambda$, $u = 0$.

\subsection{Step 4: Two-level inference falls out}
\label{ssec:twolevel}

\noindent\textit{How is the same update fast for $s_t$ and slow for
$\theory$?} A second mechanism is not needed. The asymmetry between
the factors is a choice of \emph{where} conservatism applies in
Eq.~\eqref{eq:motiv}, not two kinds of inference.

The world state sits at the ordinary-Bayes corner. Its prior is the
one-step prediction
$\sum_{s_{t-1}} B_s(s_t \mid s_{t-1}, a_{t-1})\, q(s_{t-1} \mid \theory)$
--- present because $B_s$ has nontrivial dynamics --- and at
$\alpha = \lambda$, $u = 0$, Eq.~\eqref{eq:motiv} collapses to the
exact Bayes filter
\begin{equation}
q(s_t \mid \theory) \;\propto\;
p(o_t \mid s_t, \theory)
\sum_{s_{t-1}} B_s(s_t \mid s_{t-1}, a_{t-1})\, q(s_{t-1} \mid \theory).
\label{eq:fast}
\end{equation}
The world state carries no conservatism because there is no standing
commitment to resist revising: the regime is expected to change at
every step, so tracking it responsively is correct, not motivated.

The paradigm sits away from the corner, with $\alpha < \lambda$ and
possibly $u(\theory) \neq 0$. Its prior is the carried-over belief
$q(\theory)$ --- no prediction step is needed since
$B_\theory \approx I$ --- and the evidence is the marginal likelihood
each paradigm assigns the new datum,
\begin{equation}
L(\theory \mid o_t)
\;=\;
\sum_{s_t} p(o_t \mid s_t, \theory)\, q(s_t \mid \theory).
\label{eq:marglik}
\end{equation}
Instantiating Eq.~\eqref{eq:motiv} gives
\begin{equation}
q'(\theory)
\;\propto\;
q(\theory)\, L(\theory \mid o_t)^{\alpha/\lambda}\,
e^{\, u(\theory)/\lambda}.
\label{eq:slow}
\end{equation}
The two timescales are two settings of one rule, not two mechanisms.

Eq.~\eqref{eq:slow} also reads $\lambda$. Because $\theory$ conditions
the likelihood of every past datum (Eq.~\ref{eq:gen}), adopting a new
paradigm means re-reading the entire accumulated record under a new
likelihood. For a single agent this is the cost of recomputing how
prior observations should now be understood; for a community it is
an \emph{institutional-memory} cost --- textbooks, curricula,
calibrated instruments, accepted exemplars
\citep{Heijdra1986_6,Wahyudi2026_2} all encode the prevailing
$\theory$, and a paradigm shift devalues that shared capital. We
read $\lambda$'s magnitude as a proxy for that retroactive work; this
gives a formal counterpart of Kuhnian incommensurability as
informational re-coding rather than as a departure from rationality.

\subsection{Step 5: Active sampling under expected free energy}
\label{ssec:efe}

\noindent\textit{How does paradigm-dependent attention enter without
a bespoke bias term?} Eq.~\eqref{eq:slow} only moves $\theory$ when
the marginal likelihood ratio is favourable, i.e.\ when the agent
\emph{samples} data with large $d(a)$ from Eq.~\eqref{eq:discrim}.
What the agent gets to see is therefore part of the model, not an
exogenous schedule. Action is selected by expected free energy
\citep{friston2015ActiveInferenceEpistemic,millidge2021EFE,Mirza2019_2},
\[
G(a) \;=\;
\underbrace{-\,\E{\log p(o \mid C)}}_{\text{pragmatic}}
+\underbrace{\E{D_{\mathrm{KL}}[q(s \mid o, a) \,\|\, q(s)]}}_{\text{epistemic}},
\]
with $C$ the agent's preferences. Two equivalent routes produce the
paradigm-dependence of sampling. First, the precision allocation on
the channels of the likelihood becomes paradigm-conditioned: raising
precision on channels the current paradigm reads well, and lowering
it on channels that distinguish the rival paradigm, reweights the
epistemic value of every action
\citep{feldmanFriston2010Attention,parrFriston2017Attention,
parrFriston2019AdvancedAtt,limanowskiFriston2018Dark}. Second, a
small affective utility $u(\theory)$ tilting the planning belief
makes discriminating experiments --- the ones with large $d(a)$ ---
appear risky to the preferred paradigm and so under-sampled
\citep{albarracin2022EpistemicCommunities}. Both routes deliver the
same observable: the agent does not run the experiments that would
move its belief.

The three Kuhnian faces of §\ref{sec:phenomenon} are now one
construction. Theory-laden perception lives in the likelihood
$p(o \mid s, \theory)$ of Step~2. Theory-laden experimentation lives
in the EFE policy of this step. Theory-laden hypothesis space lives
in the $\lambda$-cost of Step~3, which makes alternatives nominally
in the agent's $\theory$ support but practically unreachable on the
relevant timescale. No additive bias terms have been introduced.

\subsection{Step 6: Trust-matrix coupling and the multi-agent step}
\label{ssec:trust}

\noindent\textit{How does the social signal enter, and why does it
look the way it does?} Paradigms persist across many agents, not
one. The agent is therefore one of $N$ on a graph
$\mathcal G = (V, E)$, each carrying its own
$(q_i(s_t \mid \theory),\, q_i(\theory),\, \lambda_i,\, U_i)$. The
question is: what does a peer's communication look like as an input
to Eq.~\eqref{eq:slow}?

The natural answer is that each agent emits its current paradigm
belief, and agent $i$ aggregates its neighbours' emissions into a
trust-weighted social observation
$o^{\mathrm s}_i = \sum_j T_{ij}\, q_j(\theory)$, where $T$ is the
row-stochastic \emph{trust matrix} obtained by row-normalising the
graph adjacency $A$,
$T = \mathrm{rownorm}(A)$. The matrix carries no self term: own-belief
inertia is already represented by the prior in
Eq.~\eqref{eq:slow}, weighted against incoming evidence by
$\lambda$, so a self-weight in $T$ would duplicate it. The trust
matrix therefore encodes only which neighbours an agent attends to,
a role distinct from $\lambda$. The social signal enters
Eq.~\eqref{eq:slow} as a second accuracy channel
\citep{albarracin2022EpistemicCommunities,veissiere2020TTOM}; its
precision is set by the channel weight $\kappa$, and a confirmation
bias $\gamma$ up-weights aligned neighbours.

\begin{figure}[t]
\centering
\begin{minipage}[b]{0.35\textwidth}\centering
\begin{tikzpicture}[scale=0.92]
\node[anode] (n0) at (0.00,2.45) {0};
\node[anode] (n1) at (-1.22,2.12) {1};
\node[anode] (n2) at (-2.12,1.23) {2};
\node[anode] (n3) at (-2.45,0.00) {3};
\node[anode] (n4) at (-2.12,-1.22) {4};
\node[anode] (n5) at (-1.22,-2.12) {5};
\node[anode] (n6) at (-0.00,-2.45) {6};
\node[anode] (n7) at (1.22,-2.12) {7};
\node[anode] (n8) at (2.12,-1.23) {8};
\node[anode] (n9) at (2.45,-0.00) {9};
\node[anode] (n10) at (2.12,1.23) {10};
\node[anode] (n11) at (1.23,2.12) {11};
\draw[chan] (n0) -- (n1);
\draw[chan] (n0) -- (n2);
\draw[chan] (n0) -- (n11);
\draw[chan] (n1) -- (n5);
\draw[chan] (n1) -- (n6);
\draw[chan] (n1) -- (n7);
\draw[chan] (n1) -- (n9);
\draw[chan] (n1) -- (n10);
\draw[chan] (n1) -- (n11);
\draw[chan] (n2) -- (n3);
\draw[chan] (n2) -- (n8);
\draw[chan] (n3) -- (n4);
\draw[chan] (n3) -- (n9);
\draw[chan] (n3) -- (n10);
\draw[chan] (n4) -- (n6);
\draw[chan] (n4) -- (n8);
\draw[chan] (n5) -- (n7);
\draw[chan] (n6) -- (n7);
\draw[chan] (n6) -- (n8);
\draw[chan] (n7) -- (n9);
\draw[chan] (n8) -- (n9);
\draw[chan] (n8) -- (n10);
\draw[chan] (n9) -- (n10);
\draw[chan] (n10) -- (n11);

\end{tikzpicture}\\[2pt]
{\footnotesize (a) agents and channels}
\end{minipage}\hfill
\begin{minipage}[b]{0.31\textwidth}\centering
\begin{tikzpicture}
\input{_trustfig_adj}
\node[font=\scriptsize] at (2.04,0.65) {$j$};
\node[font=\scriptsize,rotate=90] at (-0.45,-1.7) {$i$};
\end{tikzpicture}\\[2pt]
{\footnotesize (b) adjacency $A$}
\end{minipage}\hfill
\begin{minipage}[b]{0.31\textwidth}\centering
\begin{tikzpicture}
\fill[trustc!67] (0.340,0.000) rectangle (0.680,0.340);
\fill[trustc!67] (0.680,0.000) rectangle (1.020,0.340);
\fill[trustc!67] (3.740,0.000) rectangle (4.080,0.340);
\fill[trustc!29] (0.000,-0.340) rectangle (0.340,0.000);
\fill[trustc!29] (1.700,-0.340) rectangle (2.040,0.000);
\fill[trustc!29] (2.040,-0.340) rectangle (2.380,0.000);
\fill[trustc!29] (2.380,-0.340) rectangle (2.720,0.000);
\fill[trustc!29] (3.060,-0.340) rectangle (3.400,0.000);
\fill[trustc!29] (3.400,-0.340) rectangle (3.740,0.000);
\fill[trustc!29] (3.740,-0.340) rectangle (4.080,0.000);
\fill[trustc!67] (0.000,-0.680) rectangle (0.340,-0.340);
\fill[trustc!67] (1.020,-0.680) rectangle (1.360,-0.340);
\fill[trustc!67] (2.720,-0.680) rectangle (3.060,-0.340);
\fill[trustc!50] (0.680,-1.020) rectangle (1.020,-0.680);
\fill[trustc!50] (1.360,-1.020) rectangle (1.700,-0.680);
\fill[trustc!50] (3.060,-1.020) rectangle (3.400,-0.680);
\fill[trustc!50] (3.400,-1.020) rectangle (3.740,-0.680);
\fill[trustc!67] (1.020,-1.360) rectangle (1.360,-1.020);
\fill[trustc!67] (2.040,-1.360) rectangle (2.380,-1.020);
\fill[trustc!67] (2.720,-1.360) rectangle (3.060,-1.020);
\fill[trustc!100] (0.340,-1.700) rectangle (0.680,-1.360);
\fill[trustc!100] (2.380,-1.700) rectangle (2.720,-1.360);
\fill[trustc!50] (0.340,-2.040) rectangle (0.680,-1.700);
\fill[trustc!50] (1.360,-2.040) rectangle (1.700,-1.700);
\fill[trustc!50] (2.380,-2.040) rectangle (2.720,-1.700);
\fill[trustc!50] (2.720,-2.040) rectangle (3.060,-1.700);
\fill[trustc!50] (0.340,-2.380) rectangle (0.680,-2.040);
\fill[trustc!50] (1.700,-2.380) rectangle (2.040,-2.040);
\fill[trustc!50] (2.040,-2.380) rectangle (2.380,-2.040);
\fill[trustc!50] (3.060,-2.380) rectangle (3.400,-2.040);
\fill[trustc!40] (0.680,-2.720) rectangle (1.020,-2.380);
\fill[trustc!40] (1.360,-2.720) rectangle (1.700,-2.380);
\fill[trustc!40] (2.040,-2.720) rectangle (2.380,-2.380);
\fill[trustc!40] (3.060,-2.720) rectangle (3.400,-2.380);
\fill[trustc!40] (3.400,-2.720) rectangle (3.740,-2.380);
\fill[trustc!40] (0.340,-3.060) rectangle (0.680,-2.720);
\fill[trustc!40] (1.020,-3.060) rectangle (1.360,-2.720);
\fill[trustc!40] (2.380,-3.060) rectangle (2.720,-2.720);
\fill[trustc!40] (2.720,-3.060) rectangle (3.060,-2.720);
\fill[trustc!40] (3.400,-3.060) rectangle (3.740,-2.720);
\fill[trustc!40] (0.340,-3.400) rectangle (0.680,-3.060);
\fill[trustc!40] (1.020,-3.400) rectangle (1.360,-3.060);
\fill[trustc!40] (2.720,-3.400) rectangle (3.060,-3.060);
\fill[trustc!40] (3.060,-3.400) rectangle (3.400,-3.060);
\fill[trustc!40] (3.740,-3.400) rectangle (4.080,-3.060);
\fill[trustc!67] (0.000,-3.740) rectangle (0.340,-3.400);
\fill[trustc!67] (0.340,-3.740) rectangle (0.680,-3.400);
\fill[trustc!67] (3.400,-3.740) rectangle (3.740,-3.400);
\draw[mframe] (0,0.340) rectangle (4.080,-3.740);

\node[font=\scriptsize] at (2.04,0.65) {$j$};
\node[font=\scriptsize,rotate=90] at (-0.45,-1.7) {$i$};
\end{tikzpicture}\\[2pt]
{\footnotesize (c) trust matrix $T$}
\end{minipage}
\caption{The trust network, from the Watts--Strogatz construction
($N{=}12$, $k{=}4$, rewiring $0.25$) used in the implementation
\citep{wattsStrogatz1998SmallWorld}.
(a)~Agents and their communication channels. (b)~Binary adjacency
matrix $A$. (c)~Row-stochastic trust matrix
$T = \mathrm{rownorm}(A)$: row $i$ gives the weights with which
agent $i$ mixes its neighbours' paradigm beliefs into its social
observation. The diagonal is empty because own-belief inertia is
supplied by $\lambda$, not by $T$.}
\label{fig:trust}
\end{figure}

The construction recovers, as derived objects, the primitives of the
network-epistemology tradition: scalar Bayesian updating on a graph
\citep{Bala1998_1}, DeGroot-style consensus dynamics, and the
Friedkin--Johnsen susceptibility coefficient as the linear limit of
the $\lambda$-modulated update \citep{Boudourides2026_7}.

\subsection{Step 7: Population dynamics and the damped oscillation}
\label{ssec:population}

\noindent\textit{What is the asymptotic behaviour of a population
coupled this way?} With one-round transmission delays through the
graph, the coupled population is a second-order system. Its approach
to consensus is a \emph{damped oscillation}, with damping ratio set
by the distribution of conservatism $\{\lambda_i\}$ and the spectral
structure of $T$. Specifically, the rate at which beliefs synchronise
is controlled by the algebraic connectivity (Fiedler value
$\lambda_2$) of the Laplacian. Figure~\ref{fig:anticipated} sketches
the regimes.

\begin{figure}[t]
\centering
\begin{minipage}{0.49\textwidth}\centering
\begin{tikzpicture}
\begin{axis}[width=\textwidth,height=5.2cm,
  xlabel={time}, ylabel={population mean $\bar q(\theory)$},
  xmin=0,xmax=40,ymin=0,ymax=1,ytick={0,0.5,1},
  legend style={font=\tiny,at={(0.98,0.04)},anchor=south east,draw=gray!40},
  domain=0:40,samples=120,thick,every axis plot/.append style={line width=0.7pt}]
\addplot[paradB]{0.7-0.4*exp(-0.035*x)*cos(deg(0.55*x))};
\addlegendentry{weakly damped (ringing)}
\addplot[paradA]{0.7-0.4*exp(-0.10*x)*cos(deg(0.55*x))};
\addlegendentry{under-damped}
\addplot[util]{0.7-0.4*exp(-0.22*x)*(1+0.22*x)};
\addlegendentry{critically damped}
\addplot[social]{0.7-0.4*exp(-0.13*x)};
\addlegendentry{over-damped}
\draw[gray,dashed] (axis cs:0,0.7)--(axis cs:40,0.7);
\end{axis}
\end{tikzpicture}\\[2pt]
{\footnotesize (a) damping regimes vs.\ conservatism distribution}
\end{minipage}\hfill
\begin{minipage}{0.49\textwidth}\centering
\begin{tikzpicture}
\begin{axis}[width=\textwidth,height=5.2cm,
  xlabel={time}, ylabel={belief in best-fit paradigm},
  xmin=0,xmax=40,ymin=0,ymax=1,ytick={0,0.5,1},
  legend style={font=\tiny,at={(0.98,0.5)},anchor=east,draw=gray!40},
  domain=0:40,samples=120,thick,every axis plot/.append style={line width=0.7pt}]
\addplot[paradA]{1-0.5*exp(-0.25*x)};
\addlegendentry{low $\lambda$: fast}
\addplot[social]{1-0.5*exp(-0.06*x)};
\addlegendentry{high $\lambda$: slow}
\addplot[util,dashed]{0.5+0.0001*x};
\addlegendentry{frozen ($\lambda$ large)}
\addplot[paradB]{0.12+0.38*exp(-0.13*x)};
\addlegendentry{affective lock-in}
\draw[gray,dashed] (axis cs:0,0.5)--(axis cs:40,0.5);
\end{axis}
\end{tikzpicture}\\[2pt]
{\footnotesize (b) individual regimes: convergence vs.\ lock-in}
\end{minipage}
\caption{Anticipated target dynamics (schematic). (a)~The
population's mean paradigm belief approaches consensus (dashed) as a
damped oscillation; the regime --- weakly damped, under-damped,
critically damped, over-damped --- is set by the distribution
$\{\lambda_i\}$ and the spectrum of $T$. (b)~At the agent level,
finite conservatism yields convergence to the best-fit paradigm
(fast or slow). Under unbounded $\lambda$ the belief stays frozen at
its prior; a supra-threshold affective utility produces lock-in on
the preferred non-best-fit paradigm.}
\label{fig:anticipated}
\end{figure}

The derived consequences the literature has accumulated as separate
results recover from this picture. The Zollman speed--accuracy
tradeoff% TODO ADD zollman2007NetworkEpistemology
is the topology-level statement of how $\lambda_2$ governs the
transient: denser graphs synchronise faster but lose the transient
diversity that protects against locking on the worse paradigm.
Paradoxical stubbornness --- the regularising effect of high-$\lambda_i$
agents on the population --- is the same fact read at the agent
level. The bubble-versus-chamber distinction \citep{Nguyen2018_1}
recovers as the formal separation between sparsity of $T$ (a bubble:
exposure dissolves it) and the precision weights $T$ carries on its
non-zero entries (a chamber: exposure reinforces it). The model does
not stipulate these phenomena; it \emph{produces} them.

% =====================================================================
\section{Open Question: External Curation of Attention}
\label{sec:discussion}
% =====================================================================

The most consequential open question is the model's empirical
domain. Throughout the construction the agent's policy over which
neighbours, papers, and observations to attend to is endogenous ---
selected by expected free energy under the agent's own generative
model and affective utility. In contemporary scientific practice
this assumption is strained: citation recommenders, search-engine
ranking, social-media amplification, and journal-level visibility
filters increasingly mediate which findings any individual scientist
ever encounters \citep{kulveit2025GradualDisempowerment,
critchKrueger2020ARCHES}. A curator $\mathcal{C}$ --- algorithmic,
institutional, or both --- intervenes between the latent paradigm
field and the per-agent observation channel, with its own objective
that is not in general aligned with predictive accuracy on the
underlying scientific question. Whether this collapses to a
renormalisation of the existing trust dynamics, or introduces
qualitatively new attractors that survive cross-surprisal because
the disconfirming observations are never sampled, is the natural
next question the framework writes.
